\section{Introduction}\label{sec:introduction}

Catalan's constant is the real number
\[
G=\beta(2)=\sum_{j=0}^{\infty}\frac{(-1)^j}{(2j+1)^2},
\qquad
\beta(s)=\sum_{j=0}^{\infty}\frac{(-1)^j}{(2j+1)^s}\quad(s>0).
\]
It is the simplest even value of the Dirichlet beta function, or
equivalently $L(2,\chi_{-4})$ for the odd quadratic character of
conductor $4$. Its arithmetic illustrates a basic difficulty in the
study of special values: rapidly converging rational approximations
need not be sufficiently accurate after their denominators are cleared.
We prove the following statement about this individual value.

\begin{theorem}\label{thm:main}
Catalan's constant is irrational.
\end{theorem}

\subsection{Earlier results and the approximation problem}

For odd positive integers, the beta values are rational multiples of the
corresponding powers of $\pi$. The even values present a different
arithmetic problem \cite[Introduction]{RivoalZudilin2003}.
Rivoal and Zudilin proved that infinitely many even beta values are
irrational, and that at least one of
$\beta(2),\beta(4),\ldots,\beta(14)$ is irrational
\cite[Theorems~1 and~2]{RivoalZudilin2003}.
Zudilin reduced this finite collection to the six values through
$\beta(12)$ \cite[Theorem~1]{Zudilin2019}.
Lai and Zhou subsequently reduced it to the five values
$\beta(2),\beta(4),\beta(6),\beta(8),\beta(10)$
\cite[Theorem~6.1]{LaiZhou2022}.
Fischler obtained stronger quantitative results for irrationality and
linear independence in families of Dirichlet $L$-values
\cite[Theorems~1 and~2]{Fischler2020}.
Such family results guarantee irrational members without identifying
the particular value $\beta(2)$.

For an individual value, the basic approximation criterion requires
nonzero integer linear forms $A_mG-B_m$ that tend to zero.
Convergence of $B_m/A_m$ to $G$ alone does not suffice: multiplication
by $A_m$ may remove the decay.
Ap\'ery's recurrence constructions for $\zeta(2)$ and $\zeta(3)$
\cite{Apery1979}, and Beukers's integral proofs of their irrationality
\cite{Beukers1979} succeed because the analytic decay survives the
necessary denominator clearing.
For Catalan's constant, Zudilin constructed Ap\'ery-like recurrences,
a continued fraction and double-integral representations
\cite[Theorems~1--3]{Zudilin2003}.
His discussion following Theorem~1 makes the obstruction explicit:
the displayed integer linear forms do not tend to zero, despite the
rapid convergence of their rational quotients.

The denominators themselves became an important part of the problem.
Rivoal connected Pad\'e approximation with the hypergeometric
construction and proved its conjectured denominator bounds
\cite[Theorems~1 and~2]{Rivoal2006Catalan}.
Krattenthaler and Rivoal subsequently gave a more direct hypergeometric
proof of these bounds \cite[Theorems~1 and~2]{KrattenthalerRivoal2008Catalan}.
These results establish arithmetic cancellation that is invisible in
a crude estimate of the individual summands. They still leave a
denominator cost too large for the small-linear-form criterion.
Krattenthaler and Zudilin later identified two apparently different
hypergeometric constructions of the same approximants
\cite[Section~2 and Theorem~2]{KrattenthalerZudilin2019}.
Such identities matter arithmetically because different expressions can
make different denominator factors visible.

Nesterenko developed effective approximations using half-integer
hypergeometric series and double Euler integrals
\cite{Nesterenko2016Catalan}.
Viola announced joint work with Marcovecchio improving an approximation
exponent through the Rhin--Viola permutation-group method
\cite{Viola2022Catalan}; their report gives the exponent $0.6293\ldots$
for its explicit approximants. A recent preprint of Eskandari constructs
further explicit rational approximations satisfying
$0<|G-p_m/q_m|\le q_m^{-0.62}$ for all sufficiently large $m$
\cite[Theorem~1.1]{Eskandari2026Approximations}.
These are bounds for particular constructions, not irrationality
measures. An error bound with exponent below $1$ does not by itself
make the integer forms $q_mG-p_m$ tend to zero.

Results at other characters and other places give useful comparisons.
Calegari, Dimitrov and Tang proved the $\mathbb Q$-linear independence
of $1$, $\pi^2$ and $L(2,\chi_{-3})$
\cite[Theorem~A]{CalegariDimitrovTang2024}.
Their character has conductor $3$, rather than the conductor $4$ of $G$.
Their discussion of two Catalan approximation families again identifies
the denominator cost as an obstruction for those constructions
\cite[Remark~11.1.17]{CalegariDimitrovTang2024}.
Calegari proved irrationality of a $2$-adic analogue
$G_2\in\mathbb Q_2$ \cite[Theorem~4.2]{Calegari2005}.
As explained by Calegari, Dimitrov and Tang
\cite[Section~5.2, footnote~2]{CalegariDimitrovTang2025}, the related
real identity contains an additional $\pi^2$ term that is absent in
the $2$-adic identity. This illustrates why the $2$-adic result does
not settle the real problem. In our construction, cancellation of an
additional $\zeta(2)$ term is likewise essential, although the moment
identities and the cancellation argument are proved directly below.

Sun~\cite[Theorem~1.1]{Sun2026} has also announced a proof of the same
qualitative irrationality statement. No result from that preprint is
used here.

\subsection{The determinant strategy}

We construct determinants $\Delta_N$ of size $n=48N$ whose entries
combine moments of two elementary kernels. Each moment is a rational
linear combination of $1$, $G$ and $\zeta(2)$. The polynomial rows are
chosen to have high Taylor contact: two prescribed expressions in the
rows have identical initial Taylor coefficients. This agreement cancels the $\zeta(2)$ term
in every entry. Consequently, the single hypothesis $G\in\mathbb Q$
makes all the determinants rational.

For a nonzero rational number, its ordinary absolute value is determined
by its prime valuations. Bounds on the denominators of $\Delta_N$
therefore give a lower bound on $\log|\Delta_N|$. An integral formula
for the same determinant gives an upper bound. The contradiction comes
from making these bounds incompatible as the size grows.
Zudilin's determinantal criterion gives a useful precedent for this
comparison \cite[Proposition~2 and Section~2]{Zudilin2017Determinantal}.
There, a positive moment representation produces Hankel determinants
with squared-Vandermonde integrals. Positivity supplies nonvanishing,
and the denominator and integral estimates are compared on the scale
of the square of the matrix size. His discussion of Catalan's constant
explains why the denominator growth of the approximation family
considered there still prevents application of the criterion
\cite[Section~6]{Zudilin2017Determinantal}.
Our determinant is mixed and signed; its nonvanishing and its
real-place estimate require separate arguments.

Three features of the construction make this comparison possible.
First, integral Chebyshev rows provide both the Taylor contact and
useful divisibility. At the large odd primes relevant to the leading
bound, denominators of order $p^2$ and $p$ must both be controlled.
The arithmetic estimate follows these two layers separately; cancellation
at the first layer alone would not give the required denominator bound.

Second, nonvanishing is arithmetic. Under the hypothesis $G\in\mathbb Q$,
when $N=p$ is itself a sufficiently large prime, Frobenius and a pair of
polynomial bases reduce the growing matrix to three fixed rational
matrices. An exact finite certificate proves their nonvanishing.
The reduction then gives nonzero determinants along an unbounded
prime sequence, without asserting positivity of the mixed determinant.

Third, the real estimate respects both branches of the rational
parametrization used to construct the rows. A bounded holomorphic
interpolant controls their mixed evaluations in one range of
configurations; a Hadamard bound treats the complementary range.
The interpolation estimate is uniform even when nodes approach one
another. Its proof uses a finite-dimensional Hardy-space operator,
so the evaluation determinant cancels algebraically rather than
introducing an inverse-Vandermonde estimate.
Andr\'eief's integration identity \cite{Forrester2018Andreief},
Cauchy's double alternant \cite{Krattenthaler1999}, and the
kernel/compression viewpoint of analytic interpolation
\cite{Sarason1967} supply the classical context; the
specific estimates are proved below.

The resulting Vandermonde products have logarithmic interactions,
as in logarithmic potential theory \cite{Saff2010}.
A Chebyshev expansion of the logarithmic kernel
\cite[Lemma~3.1]{GaroufalidisPopescu2013} reduces the estimate to
quadratic sums and two one-variable functions in each case. We regularize
their interactions together, control the omitted diagonals and endpoint
terms, and only then take a supremum and pass to the limit.
Explicit rational trial coefficients finish the argument. Their
certificate exhausts all stationary points and bounds entire root
brackets, rather than relying on a numerical search for the maxima.

For orientation, the two estimates are stated for the same normalized
logarithm
\begin{equation}\label{eq:normalized-logarithm}
\mathcal L_N=\frac{\log|\Delta_N|}{(48N)^2}-\frac12\log 2.
\end{equation}
Under the rationality hypothesis, Proposition~\ref{prop:finite-lower}
gives $\liminf\mathcal L_N>-2.29084$ along every sequence of nonzero
determinants, and Proposition~\ref{prop:nonvanishing} supplies such a
sequence with $N$ prime. Independently,
Proposition~\ref{prop:real-upper} gives
$\limsup\mathcal L_N\le-2.290939875<-2.2909$.
These inequalities are incompatible. The argument proves qualitative
irrationality; a bound for an irrationality measure would require
additional control of rational approximations.
Section~\ref{sec:conclusion} also deduces irrationality of the minimum
volume of an orientable complete finite-volume hyperbolic three-manifold
with exactly two cusps, and of certain arithmetic hyperbolic volumes.

\subsection{Organization and conventions}

Section~\ref{sec:foundations} constructs the rows and moments, proves
cancellation and rationality, and gives the estimate at the prime $2$.
Section~\ref{sec:odd-primes} treats odd primes and derives the
finite-place lower bound. Its only prime-distribution input is the
ordinary prime number theorem; the weighted consequence needed here
is proved locally. Section~\ref{sec:nonvanishing} establishes
nonvanishing at prime scales. Sections~\ref{sec:real-place}
and~\ref{sec:energy} give the uniform real-place and energy estimates.
Section~\ref{sec:certificate} supplies the rational data certifying the real-place bound,
and Section~\ref{sec:conclusion} assembles the contradiction.

All logarithms are natural. We use $v_p(p)=1$ and $v_p(0)=+\infty$;
for real estimates we put $\log 0=-\infty$.
Constants in asymptotic estimates may depend on the fixed rational value
hypothetically assigned to $G$, but never on $N$ or on the integration
points. The finite certificates are specified by rational data and
arithmetic instructions in the text; supplementary programs reproduce them.
